% Additional structural scope of the one-row insertion model, independent
% of any Hamilton-cycle discovery algorithm.
\paragraph{Two further reserves are incompatible with this insertion.}
The new-row leaves force matching cuts $12303\,12323$ and
$13203\,23203$, and the new-row components at these endpoints are the
two leaf $K_2$'s. In particular $12323$ cannot have another new-row edge.
In the old second-row cycle, $12023$ has only its protected edge
$12023\,13023$, and $12320$ has its protected edge $12320\,32320$
and the edge to the already matched $12323$.
Neither $12023$ nor $12320$ can therefore be a matching endpoint.
Their new-row degrees must both be two, forcing the edges
$12020\,12023$ and $12020\,12320$. Thus $12020$ is full, and the old
reserve $12020\,12030$ cannot be retained.

If the reserve $13020\,13030$ is retained, its two endpoints already
have their protected edges $13020\,13023$ and $13030\,13230$.
Hence $12030$ can meet neither $12020$ nor $13030$ in the new row.
Its only remaining edge is to $32030$, making $12030$ a matching
endpoint. Its old second-row neighbours are $12020$ and $32030$;
the former is full, so the matching must use $12030\,32030$.
But then $32030$ is also a new-row endpoint and must retain its
protected edge $32030\,32031$, leaving no room for $12030\,32030$.
This contradiction also excludes $13020\,13030$ in the general
matching-and-linear-forest insertion model, not only its three-cut case.
